
\documentclass{article}
\usepackage{colortbl}
\usepackage{makecell}
\usepackage{multirow}
\usepackage{supertabular}

\begin{document}

\newcounter{utterance}

\centering \large Interaction Transcript for game `dond', experiment `semi\_en', episode 17 with Qwen3{-}32B{-}t0.0.
\vspace{24pt}

{ \footnotesize  \setcounter{utterance}{1}
\setlength{\tabcolsep}{0pt}
\begin{supertabular}{c@{$\;$}|p{.15\linewidth}@{}p{.15\linewidth}p{.15\linewidth}p{.15\linewidth}p{.15\linewidth}p{.15\linewidth}}
    \# & $\;$A & \multicolumn{4}{c}{Game Master} & $\;\:$B\\
    \hline

    \theutterance \stepcounter{utterance}  
    & & \multicolumn{4}{p{0.6\linewidth}}{
        \cellcolor[rgb]{0.9,0.9,0.9}{
            \makecell[{{p{\linewidth}}}]{
                \texttt{\tiny{[P1$\langle$GM]}}
                \texttt{You are playing a negotiation game in which you have to agree on how to divide a set of items among you and another player.} \\
\\ 
\texttt{Rules:} \\
\texttt{(a) You and the other player are given a set of items. Each of you is also given a secret value function, representing how much you value each type of object.} \\
\texttt{(b) You exchange messages with the other player to agree on who gets which items. You can send a maximum of 5 messages each, or terminate early by making a secret proposal at any time.} \\
\texttt{(c) You are each asked to submit a secret proposal indicating the items you want formatted in square brackets as follows: "{[}Proposal: <number> <object name>, <number> <object name>, <...>{]}"} \\
\texttt{(d) If your proposals are complementary, i.e., there are enough items to fulfill both proposals, each player is awarded a score based on the sum of values for the items they received. Otherwise, both of you get zero points.} \\
\texttt{(e) Your goal is to maximize the score you receive.} \\
\\ 
\texttt{Let us start.} \\
\\ 
\texttt{The set of available items is:} \\
\\ 
\texttt{2 plates, 1 coin, 2 candies, 2 lemons} \\
\\ 
\texttt{Your secret value function is:} \\
\\ 
\texttt{plate: 0, coin: 2, candy: 4, lemon: 0} \\
\\ 
\texttt{IMPORTANT: Your messages, unless it is the secret proposal, are directly transmitted to the other player, so do not include any response to the rules or text announcing your message. To make a secret proposal, use the indicated format. Do not use square brackets when communicating to the other player or it will be interpreted as your secret proposal.} \\
            }
        }
    }
    & \\ \\

    \theutterance \stepcounter{utterance}  
    & \multicolumn{4}{p{0.6\linewidth}}{
        \cellcolor[rgb]{0.9,0.9,0.9}{
            \makecell[{{p{\linewidth}}}]{
                \texttt{\tiny{[P1$\rangle$GM]}}
                \texttt{<think>} \\
\texttt{Okay, let's see. I need to negotiate the division of these items: 2 plates, 1 coin, 2 candies, 2 lemons. My value function is plate: 0, coin: 2, candy: 4, lemon: 0. So I really value candies the most, then the coin. The other player's value function is secret, so I don't know what they want. But I need to figure out a way to get the best possible score for myself.} \\
\\ 
\texttt{First, I should consider what items I want. Since candies are worth 4 each, getting two candies would give me 8 points. The coin is worth 2, so maybe getting that as well. But there's only one coin. The plates and lemons are worth 0 to me, so I don't want them. The other player might value those differently. For example, maybe they value plates or lemons more. } \\
\\ 
\texttt{Since the proposals need to be complementary, I need to make sure that if I ask for certain items, the other player can get the rest without overlapping. Let's think about possible scenarios. If I take both candies, that's 2 candies. Then the remaining items are 2 plates, 1 coin, 2 lemons. The other player might want the coin or maybe the plates or lemons. If they want the coin, they can take it, and the rest can be split. But if they also want candies, but there are none left, that would conflict. But since I'm taking both candies, they can't take any, so their proposal must not include candies. } \\
\\ 
\texttt{Alternatively, maybe I should propose taking one candy and the coin. That would give me 4 + 2 = 6 points. Then the other player could take the other candy and maybe other items. But if they also want the coin, that's a conflict. Hmm. } \\
\\ 
\texttt{Wait, the key is to make a proposal that the other player can accept without overlapping. So if I take the two candies (total 8 points), then the remaining items are 2 plates, 1 coin, 2 lemons. The other player could take the coin (2 points) and maybe some plates or lemons. Since their value function is secret, maybe they value the coin or other items. If they take the coin, that's fine. If they take plates or lemons, which I don't care} \\
            }
        }
    }
    & & \\ \\

    \theutterance \stepcounter{utterance}  
    & & & \multicolumn{2}{p{0.3\linewidth}}{
        \cellcolor[rgb]{0.9,0.9,0.9}{
            \makecell[{{p{\linewidth}}}]{
                \texttt{\tiny{[GM$|$GM]}}
                \texttt{continue} \\
            }
        }
    }
    & & \\ \\

    \theutterance \stepcounter{utterance}  
    & & & \multicolumn{2}{p{0.3\linewidth}}{
        \cellcolor[rgb]{0.9,0.9,0.9}{
            \makecell[{{p{\linewidth}}}]{
                \texttt{\tiny{[GM$|$GM]}}
                \texttt{<think>} \\
\texttt{Okay, let's see. I need to negotiate the division of these items: 2 plates, 1 coin, 2 candies, 2 lemons. My value function is plate: 0, coin: 2, candy: 4, lemon: 0. So I really value candies the most, then the coin. The other player's value function is secret, so I don't know what they want. But I need to figure out a way to get the best possible score for myself.} \\
\\ 
\texttt{First, I should consider what items I want. Since candies are worth 4 each, getting two candies would give me 8 points. The coin is worth 2, so maybe getting that as well. But there's only one coin. The plates and lemons are worth 0 to me, so I don't want them. The other player might value those differently. For example, maybe they value plates or lemons more. } \\
\\ 
\texttt{Since the proposals need to be complementary, I need to make sure that if I ask for certain items, the other player can get the rest without overlapping. Let's think about possible scenarios. If I take both candies, that's 2 candies. Then the remaining items are 2 plates, 1 coin, 2 lemons. The other player might want the coin or maybe the plates or lemons. If they want the coin, they can take it, and the rest can be split. But if they also want candies, but there are none left, that would conflict. But since I'm taking both candies, they can't take any, so their proposal must not include candies. } \\
\\ 
\texttt{Alternatively, maybe I should propose taking one candy and the coin. That would give me 4 + 2 = 6 points. Then the other player could take the other candy and maybe other items. But if they also want the coin, that's a conflict. Hmm. } \\
\\ 
\texttt{Wait, the key is to make a proposal that the other player can accept without overlapping. So if I take the two candies (total 8 points), then the remaining items are 2 plates, 1 coin, 2 lemons. The other player could take the coin (2 points) and maybe some plates or lemons. Since their value function is secret, maybe they value the coin or other items. If they take the coin, that's fine. If they take plates or lemons, which I don't care} \\
            }
        }
    }
    & & \\ \\

    \theutterance \stepcounter{utterance}  
    & & \multicolumn{4}{p{0.6\linewidth}}{
        \cellcolor[rgb]{0.9,0.9,0.9}{
            \makecell[{{p{\linewidth}}}]{
                \texttt{\tiny{[GM$\rangle$P2]}}
                \texttt{You are playing a negotiation game in which you have to agree on how to divide a set of items among you and another player.} \\
\\ 
\texttt{Rules:} \\
\texttt{(a) You and the other player are given a set of items. Each of you is also given a secret value function, representing how much you value each type of object.} \\
\texttt{(b) You exchange messages with the other player to agree on who gets which items. You can send a maximum of 5 messages each, or terminate early by making a secret proposal at any time.} \\
\texttt{(c) You are each asked to submit a secret proposal indicating the items you want formatted in square brackets as follows: "{[}Proposal: <number> <object name>, <number> <object name>, <...>{]}"} \\
\texttt{(d) If your proposals are complementary, i.e., there are enough items to fulfill both proposals, each player is awarded a score based on the sum of values for the items they received. Otherwise, both of you get zero points.} \\
\texttt{(e) Your goal is to maximize the score you receive.} \\
\\ 
\texttt{Let us start.} \\
\\ 
\texttt{The set of available items is:} \\
\\ 
\texttt{2 plates, 1 coin, 2 candies, 2 lemons} \\
\\ 
\texttt{Your secret value function is:} \\
\\ 
\texttt{plate: 1, coin: 4, candy: 0, lemon: 2} \\
\\ 
\texttt{IMPORTANT: Your messages, unless it is the secret proposal, are directly transmitted to the other player, so do not include any response to the rules or text announcing your message. To make a secret proposal, use the indicated format. Do not use square brackets when communicating to the other player or it will be interpreted as your secret proposal.} \\
\\ 
\\ 
\texttt{<think>} \\
\texttt{Okay, let's see. I need to negotiate the division of these items: 2 plates, 1 coin, 2 candies, 2 lemons. My value function is plate: 0, coin: 2, candy: 4, lemon: 0. So I really value candies the most, then the coin. The other player's value function is secret, so I don't know what they want. But I need to figure out a way to get the best possible score for myself.} \\
\\ 
\texttt{First, I should consider what items I want. Since candies are worth 4 each, getting two candies would give me 8 points. The coin is worth 2, so maybe getting that as well. But there's only one coin. The plates and lemons are worth 0 to me, so I don't want them. The other player might value those differently. For example, maybe they value plates or lemons more. } \\
\\ 
\texttt{Since the proposals need to be complementary, I need to make sure that if I ask for certain items, the other player can get the rest without overlapping. Let's think about possible scenarios. If I take both candies, that's 2 candies. Then the remaining items are 2 plates, 1 coin, 2 lemons. The other player might want the coin or maybe the plates or lemons. If they want the coin, they can take it, and the rest can be split. But if they also want candies, but there are none left, that would conflict. But since I'm taking both candies, they can't take any, so their proposal must not include candies. } \\
\\ 
\texttt{Alternatively, maybe I should propose taking one candy and the coin. That would give me 4 + 2 = 6 points. Then the other player could take the other candy and maybe other items. But if they also want the coin, that's a conflict. Hmm. } \\
\\ 
\texttt{Wait, the key is to make a proposal that the other player can accept without overlapping. So if I take the two candies (total 8 points), then the remaining items are 2 plates, 1 coin, 2 lemons. The other player could take the coin (2 points) and maybe some plates or lemons. Since their value function is secret, maybe they value the coin or other items. If they take the coin, that's fine. If they take plates or lemons, which I don't care} \\
            }
        }
    }
    & \\ \\

    \theutterance \stepcounter{utterance}  
    & & & \multicolumn{4}{p{0.6\linewidth}}{
        \cellcolor[rgb]{0.9,0.9,0.9}{
            \makecell[{{p{\linewidth}}}]{
                \texttt{\tiny{[GM$\langle$P2]}}
                \texttt{<think>} \\
\texttt{Okay, let's start by analyzing the problem. I need to divide the items between myself and another player, with my value function being plate: 0, coin: 2, candy: 4, lemon: 0. The other player's values are secret, but I can infer based on the items they might want. The goal is to maximize my score, which depends on getting the items I value most without conflicting with the other player's needs.} \\
\\ 
\texttt{First, I should prioritize the items I value the most. Candies are the highest value for me (4 each), so getting both candies would give me 8 points. The coin is next (2 points), so if possible, I'd like that too. However, there's only one coin. The plates and lemons are worth 0 to me, so I don't want them unless necessary.} \\
\\ 
\texttt{The challenge is to make a proposal that the other player can accept without overlapping. If I take both candies and the coin, that's 2 candies and 1 coin. The remaining items would be 2 plates and 2 lemons. The other player might value the coin, but since I took it, they can't. However, if their value function includes plates or lemons, they might be okay with getting those. For example, if they value plates or lemons, they could take those and leave the rest. But since their values are secret, I need to make a proposal that's likely to be complementary.} \\
\\ 
\texttt{Alternatively, if I take both candies (8 points) and leave the coin, plates, and lemons for the other player, that might work. The other player could take the coin (if they value it) and the plates/lemons. Since I don't care about those, it's a win for me. The key is to ensure that the other player's proposal doesn't require items I've already taken. } \\
\\ 
\texttt{Another angle: if the other player also wants candies, but there are only two, we can't both get any. But since I value candies the most, it's better for me to take both. The other player might have to settle for other items. } \\
\\ 
\texttt{So, my optimal proposal would be to take both candies and the coin. That gives me 4*2 + 2 = 10 points. The remaining items are 2 plates and 2 lemons. If the other player's proposal is to take those, then it} \\
            }
        }
    }
     \\ \\

    \theutterance \stepcounter{utterance}  
    & & & \multicolumn{2}{p{0.3\linewidth}}{
        \cellcolor[rgb]{0.9,0.9,0.9}{
            \makecell[{{p{\linewidth}}}]{
                \texttt{\tiny{[GM$|$GM]}}
                \texttt{continue} \\
            }
        }
    }
    & & \\ \\

    \theutterance \stepcounter{utterance}  
    & & & \multicolumn{2}{p{0.3\linewidth}}{
        \cellcolor[rgb]{0.9,0.9,0.9}{
            \makecell[{{p{\linewidth}}}]{
                \texttt{\tiny{[GM$|$GM]}}
                \texttt{<think>} \\
\texttt{Okay, let's start by analyzing the problem. I need to divide the items between myself and another player, with my value function being plate: 0, coin: 2, candy: 4, lemon: 0. The other player's values are secret, but I can infer based on the items they might want. The goal is to maximize my score, which depends on getting the items I value most without conflicting with the other player's needs.} \\
\\ 
\texttt{First, I should prioritize the items I value the most. Candies are the highest value for me (4 each), so getting both candies would give me 8 points. The coin is next (2 points), so if possible, I'd like that too. However, there's only one coin. The plates and lemons are worth 0 to me, so I don't want them unless necessary.} \\
\\ 
\texttt{The challenge is to make a proposal that the other player can accept without overlapping. If I take both candies and the coin, that's 2 candies and 1 coin. The remaining items would be 2 plates and 2 lemons. The other player might value the coin, but since I took it, they can't. However, if their value function includes plates or lemons, they might be okay with getting those. For example, if they value plates or lemons, they could take those and leave the rest. But since their values are secret, I need to make a proposal that's likely to be complementary.} \\
\\ 
\texttt{Alternatively, if I take both candies (8 points) and leave the coin, plates, and lemons for the other player, that might work. The other player could take the coin (if they value it) and the plates/lemons. Since I don't care about those, it's a win for me. The key is to ensure that the other player's proposal doesn't require items I've already taken. } \\
\\ 
\texttt{Another angle: if the other player also wants candies, but there are only two, we can't both get any. But since I value candies the most, it's better for me to take both. The other player might have to settle for other items. } \\
\\ 
\texttt{So, my optimal proposal would be to take both candies and the coin. That gives me 4*2 + 2 = 10 points. The remaining items are 2 plates and 2 lemons. If the other player's proposal is to take those, then it} \\
            }
        }
    }
    & & \\ \\

    \theutterance \stepcounter{utterance}  
    & & \multicolumn{4}{p{0.6\linewidth}}{
        \cellcolor[rgb]{0.9,0.9,0.9}{
            \makecell[{{p{\linewidth}}}]{
                \texttt{\tiny{[P1$\langle$GM]}}
                \texttt{<think>} \\
\texttt{Okay, let's start by analyzing the problem. I need to divide the items between myself and another player, with my value function being plate: 0, coin: 2, candy: 4, lemon: 0. The other player's values are secret, but I can infer based on the items they might want. The goal is to maximize my score, which depends on getting the items I value most without conflicting with the other player's needs.} \\
\\ 
\texttt{First, I should prioritize the items I value the most. Candies are the highest value for me (4 each), so getting both candies would give me 8 points. The coin is next (2 points), so if possible, I'd like that too. However, there's only one coin. The plates and lemons are worth 0 to me, so I don't want them unless necessary.} \\
\\ 
\texttt{The challenge is to make a proposal that the other player can accept without overlapping. If I take both candies and the coin, that's 2 candies and 1 coin. The remaining items would be 2 plates and 2 lemons. The other player might value the coin, but since I took it, they can't. However, if their value function includes plates or lemons, they might be okay with getting those. For example, if they value plates or lemons, they could take those and leave the rest. But since their values are secret, I need to make a proposal that's likely to be complementary.} \\
\\ 
\texttt{Alternatively, if I take both candies (8 points) and leave the coin, plates, and lemons for the other player, that might work. The other player could take the coin (if they value it) and the plates/lemons. Since I don't care about those, it's a win for me. The key is to ensure that the other player's proposal doesn't require items I've already taken. } \\
\\ 
\texttt{Another angle: if the other player also wants candies, but there are only two, we can't both get any. But since I value candies the most, it's better for me to take both. The other player might have to settle for other items. } \\
\\ 
\texttt{So, my optimal proposal would be to take both candies and the coin. That gives me 4*2 + 2 = 10 points. The remaining items are 2 plates and 2 lemons. If the other player's proposal is to take those, then it} \\
            }
        }
    }
    & \\ \\

    \theutterance \stepcounter{utterance}  
    & \multicolumn{4}{p{0.6\linewidth}}{
        \cellcolor[rgb]{0.9,0.9,0.9}{
            \makecell[{{p{\linewidth}}}]{
                \texttt{\tiny{[P1$\rangle$GM]}}
                \texttt{<think>} \\
\texttt{Okay, let's try to figure out the best way to approach this. My value function is plate: 0, coin: 2, candy: 4, lemon: 0. So the candies are the most valuable to me, followed by the coin. The other player's values are secret, so I don't know what they want. But I need to make sure that our proposals are complementary so we both get points.} \\
\\ 
\texttt{First, I should aim to get the items I value the most. That would be the two candies (4 each, so 8 points) and the coin (2 points). If I take both candies and the coin, that's 10 points. The remaining items would be 2 plates and 2 lemons. The other player might want those, but since I don't value them, it's okay if they take them. But I need to check if the other player's proposal would conflict with mine. If they also want the coin or candies, there's a problem. But since I'm taking the candies and coin, they can't take those. So their proposal should only include the plates and lemons. } \\
\\ 
\texttt{Wait, but the other player's value function is secret. Maybe they value the plates or lemons more. If they take the plates and lemons, that's fine. But if they also want the coin, then our proposals would conflict. But since I'm already taking the coin, they can't take it. So maybe the other player's proposal would be to take the plates and lemons. That way, our proposals are complementary. } \\
\\ 
\texttt{Alternatively, maybe I should leave the coin for the other player. If I take both candies (8 points) and leave the coin, plates, and lemons for them. Then their proposal could include the coin and maybe some plates or lemons. But if they take the coin, that's 2 points for them, and I have 8. That's a good split. } \\
\\ 
\texttt{But I need to make sure that the other player's proposal doesn't require items I've already taken. So if I take the candies and the coin, they can't take those. If they take the remaining items (plates and lemons), that's okay. But if their value function includes the coin, they might want it, but since I took it, they can't. So maybe the other player would be okay with taking the plates and lemons if they} \\
            }
        }
    }
    & & \\ \\

    \theutterance \stepcounter{utterance}  
    & & & \multicolumn{2}{p{0.3\linewidth}}{
        \cellcolor[rgb]{0.9,0.9,0.9}{
            \makecell[{{p{\linewidth}}}]{
                \texttt{\tiny{[GM$|$GM]}}
                \texttt{continue} \\
            }
        }
    }
    & & \\ \\

    \theutterance \stepcounter{utterance}  
    & & & \multicolumn{2}{p{0.3\linewidth}}{
        \cellcolor[rgb]{0.9,0.9,0.9}{
            \makecell[{{p{\linewidth}}}]{
                \texttt{\tiny{[GM$|$GM]}}
                \texttt{<think>} \\
\texttt{Okay, let's try to figure out the best way to approach this. My value function is plate: 0, coin: 2, candy: 4, lemon: 0. So the candies are the most valuable to me, followed by the coin. The other player's values are secret, so I don't know what they want. But I need to make sure that our proposals are complementary so we both get points.} \\
\\ 
\texttt{First, I should aim to get the items I value the most. That would be the two candies (4 each, so 8 points) and the coin (2 points). If I take both candies and the coin, that's 10 points. The remaining items would be 2 plates and 2 lemons. The other player might want those, but since I don't value them, it's okay if they take them. But I need to check if the other player's proposal would conflict with mine. If they also want the coin or candies, there's a problem. But since I'm taking the candies and coin, they can't take those. So their proposal should only include the plates and lemons. } \\
\\ 
\texttt{Wait, but the other player's value function is secret. Maybe they value the plates or lemons more. If they take the plates and lemons, that's fine. But if they also want the coin, then our proposals would conflict. But since I'm already taking the coin, they can't take it. So maybe the other player's proposal would be to take the plates and lemons. That way, our proposals are complementary. } \\
\\ 
\texttt{Alternatively, maybe I should leave the coin for the other player. If I take both candies (8 points) and leave the coin, plates, and lemons for them. Then their proposal could include the coin and maybe some plates or lemons. But if they take the coin, that's 2 points for them, and I have 8. That's a good split. } \\
\\ 
\texttt{But I need to make sure that the other player's proposal doesn't require items I've already taken. So if I take the candies and the coin, they can't take those. If they take the remaining items (plates and lemons), that's okay. But if their value function includes the coin, they might want it, but since I took it, they can't. So maybe the other player would be okay with taking the plates and lemons if they} \\
            }
        }
    }
    & & \\ \\

    \theutterance \stepcounter{utterance}  
    & & \multicolumn{4}{p{0.6\linewidth}}{
        \cellcolor[rgb]{0.9,0.9,0.9}{
            \makecell[{{p{\linewidth}}}]{
                \texttt{\tiny{[GM$\rangle$P2]}}
                \texttt{<think>} \\
\texttt{Okay, let's try to figure out the best way to approach this. My value function is plate: 0, coin: 2, candy: 4, lemon: 0. So the candies are the most valuable to me, followed by the coin. The other player's values are secret, so I don't know what they want. But I need to make sure that our proposals are complementary so we both get points.} \\
\\ 
\texttt{First, I should aim to get the items I value the most. That would be the two candies (4 each, so 8 points) and the coin (2 points). If I take both candies and the coin, that's 10 points. The remaining items would be 2 plates and 2 lemons. The other player might want those, but since I don't value them, it's okay if they take them. But I need to check if the other player's proposal would conflict with mine. If they also want the coin or candies, there's a problem. But since I'm taking the candies and coin, they can't take those. So their proposal should only include the plates and lemons. } \\
\\ 
\texttt{Wait, but the other player's value function is secret. Maybe they value the plates or lemons more. If they take the plates and lemons, that's fine. But if they also want the coin, then our proposals would conflict. But since I'm already taking the coin, they can't take it. So maybe the other player's proposal would be to take the plates and lemons. That way, our proposals are complementary. } \\
\\ 
\texttt{Alternatively, maybe I should leave the coin for the other player. If I take both candies (8 points) and leave the coin, plates, and lemons for them. Then their proposal could include the coin and maybe some plates or lemons. But if they take the coin, that's 2 points for them, and I have 8. That's a good split. } \\
\\ 
\texttt{But I need to make sure that the other player's proposal doesn't require items I've already taken. So if I take the candies and the coin, they can't take those. If they take the remaining items (plates and lemons), that's okay. But if their value function includes the coin, they might want it, but since I took it, they can't. So maybe the other player would be okay with taking the plates and lemons if they} \\
            }
        }
    }
    & \\ \\

    \theutterance \stepcounter{utterance}  
    & & & \multicolumn{4}{p{0.6\linewidth}}{
        \cellcolor[rgb]{0.9,0.9,0.9}{
            \makecell[{{p{\linewidth}}}]{
                \texttt{\tiny{[GM$\langle$P2]}}
                \texttt{<think>} \\
\texttt{Okay, let's try to figure out the best way to approach this. My value function is plate: 0, coin: 2, candy: 4, lemon: 0. So the candies are the most valuable to me, followed by the coin. The other player's values are secret, so I don't know what they want. But I need to make sure that our proposals are complementary so we both get points.} \\
\\ 
\texttt{First, I should aim to get the items I value the most. That would be the two candies (4 each, so 8 points) and the coin (2 points). If I take both candies and the coin, that's 10 points. The remaining items would be 2 plates and 2 lemons. The other player might want those, but since I don't value them, it's okay if they take them. But I need to check if the other player's proposal would conflict with mine. If they also want the coin or candies, there's a problem. But since I'm taking the candies and coin, they can't take those. So their proposal should only include the plates and lemons. } \\
\\ 
\texttt{Wait, but the other player's value function is secret. Maybe they value the plates or lemons more. If they take the plates and lemons, that's fine. But if they also want the coin, then our proposals would conflict. But since I'm already taking the coin, they can't take it. So maybe the other player's proposal would be to take the plates and lemons. That way, our proposals are complementary. } \\
\\ 
\texttt{Alternatively, maybe I should leave the coin for the other player. If I take both candies (8 points) and leave the coin, plates, and lemons for them. Then their proposal could include the coin and maybe some plates or lemons. But if they take the coin, that's 2 points for them, and I have 8. That's a good split. } \\
\\ 
\texttt{But I need to make sure that the other player's proposal doesn't require items I've already taken. So if I take the candies and the coin, they can't take those. If they take the remaining items (plates and lemons), that's okay. But if their value function includes the coin, they might want it, but since I took it, they can't. So maybe the other player would be okay with taking the plates and lemons if they} \\
            }
        }
    }
     \\ \\

    \theutterance \stepcounter{utterance}  
    & & & \multicolumn{2}{p{0.3\linewidth}}{
        \cellcolor[rgb]{0.9,0.9,0.9}{
            \makecell[{{p{\linewidth}}}]{
                \texttt{\tiny{[GM$|$GM]}}
                \texttt{continue} \\
            }
        }
    }
    & & \\ \\

    \theutterance \stepcounter{utterance}  
    & & & \multicolumn{2}{p{0.3\linewidth}}{
        \cellcolor[rgb]{0.9,0.9,0.9}{
            \makecell[{{p{\linewidth}}}]{
                \texttt{\tiny{[GM$|$GM]}}
                \texttt{<think>} \\
\texttt{Okay, let's try to figure out the best way to approach this. My value function is plate: 0, coin: 2, candy: 4, lemon: 0. So the candies are the most valuable to me, followed by the coin. The other player's values are secret, so I don't know what they want. But I need to make sure that our proposals are complementary so we both get points.} \\
\\ 
\texttt{First, I should aim to get the items I value the most. That would be the two candies (4 each, so 8 points) and the coin (2 points). If I take both candies and the coin, that's 10 points. The remaining items would be 2 plates and 2 lemons. The other player might want those, but since I don't value them, it's okay if they take them. But I need to check if the other player's proposal would conflict with mine. If they also want the coin or candies, there's a problem. But since I'm taking the candies and coin, they can't take those. So their proposal should only include the plates and lemons. } \\
\\ 
\texttt{Wait, but the other player's value function is secret. Maybe they value the plates or lemons more. If they take the plates and lemons, that's fine. But if they also want the coin, then our proposals would conflict. But since I'm already taking the coin, they can't take it. So maybe the other player's proposal would be to take the plates and lemons. That way, our proposals are complementary. } \\
\\ 
\texttt{Alternatively, maybe I should leave the coin for the other player. If I take both candies (8 points) and leave the coin, plates, and lemons for them. Then their proposal could include the coin and maybe some plates or lemons. But if they take the coin, that's 2 points for them, and I have 8. That's a good split. } \\
\\ 
\texttt{But I need to make sure that the other player's proposal doesn't require items I've already taken. So if I take the candies and the coin, they can't take those. If they take the remaining items (plates and lemons), that's okay. But if their value function includes the coin, they might want it, but since I took it, they can't. So maybe the other player would be okay with taking the plates and lemons if they} \\
            }
        }
    }
    & & \\ \\

    \theutterance \stepcounter{utterance}  
    & & \multicolumn{4}{p{0.6\linewidth}}{
        \cellcolor[rgb]{0.9,0.9,0.9}{
            \makecell[{{p{\linewidth}}}]{
                \texttt{\tiny{[P1$\langle$GM]}}
                \texttt{<think>} \\
\texttt{Okay, let's try to figure out the best way to approach this. My value function is plate: 0, coin: 2, candy: 4, lemon: 0. So the candies are the most valuable to me, followed by the coin. The other player's values are secret, so I don't know what they want. But I need to make sure that our proposals are complementary so we both get points.} \\
\\ 
\texttt{First, I should aim to get the items I value the most. That would be the two candies (4 each, so 8 points) and the coin (2 points). If I take both candies and the coin, that's 10 points. The remaining items would be 2 plates and 2 lemons. The other player might want those, but since I don't value them, it's okay if they take them. But I need to check if the other player's proposal would conflict with mine. If they also want the coin or candies, there's a problem. But since I'm taking the candies and coin, they can't take those. So their proposal should only include the plates and lemons. } \\
\\ 
\texttt{Wait, but the other player's value function is secret. Maybe they value the plates or lemons more. If they take the plates and lemons, that's fine. But if they also want the coin, then our proposals would conflict. But since I'm already taking the coin, they can't take it. So maybe the other player's proposal would be to take the plates and lemons. That way, our proposals are complementary. } \\
\\ 
\texttt{Alternatively, maybe I should leave the coin for the other player. If I take both candies (8 points) and leave the coin, plates, and lemons for them. Then their proposal could include the coin and maybe some plates or lemons. But if they take the coin, that's 2 points for them, and I have 8. That's a good split. } \\
\\ 
\texttt{But I need to make sure that the other player's proposal doesn't require items I've already taken. So if I take the candies and the coin, they can't take those. If they take the remaining items (plates and lemons), that's okay. But if their value function includes the coin, they might want it, but since I took it, they can't. So maybe the other player would be okay with taking the plates and lemons if they} \\
            }
        }
    }
    & \\ \\

    \theutterance \stepcounter{utterance}  
    & \multicolumn{4}{p{0.6\linewidth}}{
        \cellcolor[rgb]{0.9,0.9,0.9}{
            \makecell[{{p{\linewidth}}}]{
                \texttt{\tiny{[P1$\rangle$GM]}}
                \texttt{<think>} \\
\texttt{Okay, let's try to figure out the best way to approach this. My value function is plate: 0, coin: 2, candy: 4, lemon: 0. So the candies are the most valuable to me, followed by the coin. The other player's values are secret, so I don't know what they want. But I need to make sure that our proposals are complementary so we both get points.} \\
\\ 
\texttt{First, I should aim to get the items I value the most. That would be the two candies (4 each, so 8 points) and the coin (2 points). If I take both candies and the coin, that's 10 points. The remaining items would be 2 plates and 2 lemons. The other player might want those, but since I don't value them, it's okay if they take them. But I need to check if the other player's proposal would conflict with mine. If they also want the coin or candies, there's a problem. But since I'm taking the candies and coin, they can't take those. So their proposal should only include the plates and lemons. } \\
\\ 
\texttt{Wait, but the other player's value function is secret. Maybe they value the plates or lemons more. If they take the plates and lemons, that's fine. But if they also want the coin, then our proposals would conflict. But since I'm already taking the coin, they can't take it. So maybe the other player's proposal would be to take the plates and lemons. That way, our proposals are complementary. } \\
\\ 
\texttt{Alternatively, maybe I should leave the coin for the other player. If I take both candies (8 points) and leave the coin, plates, and lemons for them. Then their proposal could include the coin and maybe some plates or lemons. But if they take the coin, that's 2 points for them, and I have 8. That's a good split. } \\
\\ 
\texttt{But I need to make sure that the other player's proposal doesn't require items I've already taken. So if I take the candies and the coin, they can't take those. If they take the remaining items (plates and lemons), that's okay. But if their value function includes the coin, they might want it, but since I took it, they can't. So maybe the other player would be okay with taking the plates and lemons if they} \\
            }
        }
    }
    & & \\ \\

    \theutterance \stepcounter{utterance}  
    & & & \multicolumn{2}{p{0.3\linewidth}}{
        \cellcolor[rgb]{0.9,0.9,0.9}{
            \makecell[{{p{\linewidth}}}]{
                \texttt{\tiny{[GM$|$GM]}}
                \texttt{continue} \\
            }
        }
    }
    & & \\ \\

    \theutterance \stepcounter{utterance}  
    & & & \multicolumn{2}{p{0.3\linewidth}}{
        \cellcolor[rgb]{0.9,0.9,0.9}{
            \makecell[{{p{\linewidth}}}]{
                \texttt{\tiny{[GM$|$GM]}}
                \texttt{<think>} \\
\texttt{Okay, let's try to figure out the best way to approach this. My value function is plate: 0, coin: 2, candy: 4, lemon: 0. So the candies are the most valuable to me, followed by the coin. The other player's values are secret, so I don't know what they want. But I need to make sure that our proposals are complementary so we both get points.} \\
\\ 
\texttt{First, I should aim to get the items I value the most. That would be the two candies (4 each, so 8 points) and the coin (2 points). If I take both candies and the coin, that's 10 points. The remaining items would be 2 plates and 2 lemons. The other player might want those, but since I don't value them, it's okay if they take them. But I need to check if the other player's proposal would conflict with mine. If they also want the coin or candies, there's a problem. But since I'm taking the candies and coin, they can't take those. So their proposal should only include the plates and lemons. } \\
\\ 
\texttt{Wait, but the other player's value function is secret. Maybe they value the plates or lemons more. If they take the plates and lemons, that's fine. But if they also want the coin, then our proposals would conflict. But since I'm already taking the coin, they can't take it. So maybe the other player's proposal would be to take the plates and lemons. That way, our proposals are complementary. } \\
\\ 
\texttt{Alternatively, maybe I should leave the coin for the other player. If I take both candies (8 points) and leave the coin, plates, and lemons for them. Then their proposal could include the coin and maybe some plates or lemons. But if they take the coin, that's 2 points for them, and I have 8. That's a good split. } \\
\\ 
\texttt{But I need to make sure that the other player's proposal doesn't require items I've already taken. So if I take the candies and the coin, they can't take those. If they take the remaining items (plates and lemons), that's okay. But if their value function includes the coin, they might want it, but since I took it, they can't. So maybe the other player would be okay with taking the plates and lemons if they} \\
            }
        }
    }
    & & \\ \\

    \theutterance \stepcounter{utterance}  
    & & \multicolumn{4}{p{0.6\linewidth}}{
        \cellcolor[rgb]{0.9,0.9,0.9}{
            \makecell[{{p{\linewidth}}}]{
                \texttt{\tiny{[GM$\rangle$P2]}}
                \texttt{<think>} \\
\texttt{Okay, let's try to figure out the best way to approach this. My value function is plate: 0, coin: 2, candy: 4, lemon: 0. So the candies are the most valuable to me, followed by the coin. The other player's values are secret, so I don't know what they want. But I need to make sure that our proposals are complementary so we both get points.} \\
\\ 
\texttt{First, I should aim to get the items I value the most. That would be the two candies (4 each, so 8 points) and the coin (2 points). If I take both candies and the coin, that's 10 points. The remaining items would be 2 plates and 2 lemons. The other player might want those, but since I don't value them, it's okay if they take them. But I need to check if the other player's proposal would conflict with mine. If they also want the coin or candies, there's a problem. But since I'm taking the candies and coin, they can't take those. So their proposal should only include the plates and lemons. } \\
\\ 
\texttt{Wait, but the other player's value function is secret. Maybe they value the plates or lemons more. If they take the plates and lemons, that's fine. But if they also want the coin, then our proposals would conflict. But since I'm already taking the coin, they can't take it. So maybe the other player's proposal would be to take the plates and lemons. That way, our proposals are complementary. } \\
\\ 
\texttt{Alternatively, maybe I should leave the coin for the other player. If I take both candies (8 points) and leave the coin, plates, and lemons for them. Then their proposal could include the coin and maybe some plates or lemons. But if they take the coin, that's 2 points for them, and I have 8. That's a good split. } \\
\\ 
\texttt{But I need to make sure that the other player's proposal doesn't require items I've already taken. So if I take the candies and the coin, they can't take those. If they take the remaining items (plates and lemons), that's okay. But if their value function includes the coin, they might want it, but since I took it, they can't. So maybe the other player would be okay with taking the plates and lemons if they} \\
            }
        }
    }
    & \\ \\

    \theutterance \stepcounter{utterance}  
    & & & \multicolumn{4}{p{0.6\linewidth}}{
        \cellcolor[rgb]{0.9,0.9,0.9}{
            \makecell[{{p{\linewidth}}}]{
                \texttt{\tiny{[GM$\langle$P2]}}
                \texttt{<think>} \\
\texttt{Okay, let's try to figure out the best way to approach this. My value function is plate: 0, coin: 2, candy: 4, lemon: 0. So the candies are the most valuable to me, followed by the coin. The other player's values are secret, so I don't know what they want. But I need to make sure that our proposals are complementary so we both get points.} \\
\\ 
\texttt{First, I should aim to get the items I value the most. That would be the two candies (4 each, so 8 points) and the coin (2 points). If I take both candies and the coin, that's 10 points. The remaining items would be 2 plates and 2 lemons. The other player might want those, but since I don't value them, it's okay if they take them. But I need to check if the other player's proposal would conflict with mine. If they also want the coin or candies, there's a problem. But since I'm taking the candies and coin, they can't take those. So their proposal should only include the plates and lemons. } \\
\\ 
\texttt{Wait, but the other player's value function is secret. Maybe they value the plates or lemons more. If they take the plates and lemons, that's fine. But if they also want the coin, then our proposals would conflict. But since I'm already taking the coin, they can't take it. So maybe the other player's proposal would be to take the plates and lemons. That way, our proposals are complementary. } \\
\\ 
\texttt{Alternatively, maybe I should leave the coin for the other player. If I take both candies (8 points) and leave the coin, plates, and lemons for them. Then their proposal could include the coin and maybe some plates or lemons. But if they take the coin, that's 2 points for them, and I have 8. That's a good split. } \\
\\ 
\texttt{But I need to make sure that the other player's proposal doesn't require items I've already taken. So if I take the candies and the coin, they can't take those. If they take the remaining items (plates and lemons), that's okay. But if their value function includes the coin, they might want it, but since I took it, they can't. So maybe the other player would be okay with taking the plates and lemons if they} \\
            }
        }
    }
     \\ \\

    \theutterance \stepcounter{utterance}  
    & & & \multicolumn{2}{p{0.3\linewidth}}{
        \cellcolor[rgb]{0.9,0.9,0.9}{
            \makecell[{{p{\linewidth}}}]{
                \texttt{\tiny{[GM$|$GM]}}
                \texttt{continue} \\
            }
        }
    }
    & & \\ \\

    \theutterance \stepcounter{utterance}  
    & & & \multicolumn{2}{p{0.3\linewidth}}{
        \cellcolor[rgb]{0.9,0.9,0.9}{
            \makecell[{{p{\linewidth}}}]{
                \texttt{\tiny{[GM$|$GM]}}
                \texttt{<think>} \\
\texttt{Okay, let's try to figure out the best way to approach this. My value function is plate: 0, coin: 2, candy: 4, lemon: 0. So the candies are the most valuable to me, followed by the coin. The other player's values are secret, so I don't know what they want. But I need to make sure that our proposals are complementary so we both get points.} \\
\\ 
\texttt{First, I should aim to get the items I value the most. That would be the two candies (4 each, so 8 points) and the coin (2 points). If I take both candies and the coin, that's 10 points. The remaining items would be 2 plates and 2 lemons. The other player might want those, but since I don't value them, it's okay if they take them. But I need to check if the other player's proposal would conflict with mine. If they also want the coin or candies, there's a problem. But since I'm taking the candies and coin, they can't take those. So their proposal should only include the plates and lemons. } \\
\\ 
\texttt{Wait, but the other player's value function is secret. Maybe they value the plates or lemons more. If they take the plates and lemons, that's fine. But if they also want the coin, then our proposals would conflict. But since I'm already taking the coin, they can't take it. So maybe the other player's proposal would be to take the plates and lemons. That way, our proposals are complementary. } \\
\\ 
\texttt{Alternatively, maybe I should leave the coin for the other player. If I take both candies (8 points) and leave the coin, plates, and lemons for them. Then their proposal could include the coin and maybe some plates or lemons. But if they take the coin, that's 2 points for them, and I have 8. That's a good split. } \\
\\ 
\texttt{But I need to make sure that the other player's proposal doesn't require items I've already taken. So if I take the candies and the coin, they can't take those. If they take the remaining items (plates and lemons), that's okay. But if their value function includes the coin, they might want it, but since I took it, they can't. So maybe the other player would be okay with taking the plates and lemons if they} \\
            }
        }
    }
    & & \\ \\

    \theutterance \stepcounter{utterance}  
    & & \multicolumn{4}{p{0.6\linewidth}}{
        \cellcolor[rgb]{0.9,0.9,0.9}{
            \makecell[{{p{\linewidth}}}]{
                \texttt{\tiny{[P1$\langle$GM]}}
                \texttt{<think>} \\
\texttt{Okay, let's try to figure out the best way to approach this. My value function is plate: 0, coin: 2, candy: 4, lemon: 0. So the candies are the most valuable to me, followed by the coin. The other player's values are secret, so I don't know what they want. But I need to make sure that our proposals are complementary so we both get points.} \\
\\ 
\texttt{First, I should aim to get the items I value the most. That would be the two candies (4 each, so 8 points) and the coin (2 points). If I take both candies and the coin, that's 10 points. The remaining items would be 2 plates and 2 lemons. The other player might want those, but since I don't value them, it's okay if they take them. But I need to check if the other player's proposal would conflict with mine. If they also want the coin or candies, there's a problem. But since I'm taking the candies and coin, they can't take those. So their proposal should only include the plates and lemons. } \\
\\ 
\texttt{Wait, but the other player's value function is secret. Maybe they value the plates or lemons more. If they take the plates and lemons, that's fine. But if they also want the coin, then our proposals would conflict. But since I'm already taking the coin, they can't take it. So maybe the other player's proposal would be to take the plates and lemons. That way, our proposals are complementary. } \\
\\ 
\texttt{Alternatively, maybe I should leave the coin for the other player. If I take both candies (8 points) and leave the coin, plates, and lemons for them. Then their proposal could include the coin and maybe some plates or lemons. But if they take the coin, that's 2 points for them, and I have 8. That's a good split. } \\
\\ 
\texttt{But I need to make sure that the other player's proposal doesn't require items I've already taken. So if I take the candies and the coin, they can't take those. If they take the remaining items (plates and lemons), that's okay. But if their value function includes the coin, they might want it, but since I took it, they can't. So maybe the other player would be okay with taking the plates and lemons if they} \\
            }
        }
    }
    & \\ \\

    \theutterance \stepcounter{utterance}  
    & \multicolumn{4}{p{0.6\linewidth}}{
        \cellcolor[rgb]{0.9,0.9,0.9}{
            \makecell[{{p{\linewidth}}}]{
                \texttt{\tiny{[P1$\rangle$GM]}}
                \texttt{<think>} \\
\texttt{Okay, let's try to figure out the best way to approach this. My value function is plate: 0, coin: 2, candy: 4, lemon: 0. So the candies are the most valuable to me, followed by the coin. The other player's values are secret, so I don't know what they want. But I need to make sure that our proposals are complementary so we both get points.} \\
\\ 
\texttt{First, I should aim to get the items I value the most. That would be the two candies (4 each, so 8 points) and the coin (2 points). If I take both candies and the coin, that's 10 points. The remaining items would be 2 plates and 2 lemons. The other player might want those, but since I don't value them, it's okay if they take them. But I need to check if the other player's proposal would conflict with mine. If they also want the coin or candies, there's a problem. But since I'm taking the candies and coin, they can't take those. So their proposal should only include the plates and lemons. } \\
\\ 
\texttt{Wait, but the other player's value function is secret. Maybe they value the plates or lemons more. If they take the plates and lemons, that's fine. But if they also want the coin, then our proposals would conflict. But since I'm already taking the coin, they can't take it. So maybe the other player's proposal would be to take the plates and lemons. That way, our proposals are complementary. } \\
\\ 
\texttt{Alternatively, maybe I should leave the coin for the other player. If I take both candies (8 points) and leave the coin, plates, and lemons for them. Then their proposal could include the coin and maybe some plates or lemons. But if they take the coin, that's 2 points for them, and I have 8. That's a good split. } \\
\\ 
\texttt{But I need to make sure that the other player's proposal doesn't require items I've already taken. So if I take the candies and the coin, they can't take those. If they take the remaining items (plates and lemons), that's okay. But if their value function includes the coin, they might want it, but since I took it, they can't. So maybe the other player would be okay with taking the plates and lemons if they} \\
            }
        }
    }
    & & \\ \\

    \theutterance \stepcounter{utterance}  
    & & & \multicolumn{2}{p{0.3\linewidth}}{
        \cellcolor[rgb]{0.9,0.9,0.9}{
            \makecell[{{p{\linewidth}}}]{
                \texttt{\tiny{[GM$|$GM]}}
                \texttt{continue} \\
            }
        }
    }
    & & \\ \\

    \theutterance \stepcounter{utterance}  
    & & & \multicolumn{2}{p{0.3\linewidth}}{
        \cellcolor[rgb]{0.9,0.9,0.9}{
            \makecell[{{p{\linewidth}}}]{
                \texttt{\tiny{[GM$|$GM]}}
                \texttt{<think>} \\
\texttt{Okay, let's try to figure out the best way to approach this. My value function is plate: 0, coin: 2, candy: 4, lemon: 0. So the candies are the most valuable to me, followed by the coin. The other player's values are secret, so I don't know what they want. But I need to make sure that our proposals are complementary so we both get points.} \\
\\ 
\texttt{First, I should aim to get the items I value the most. That would be the two candies (4 each, so 8 points) and the coin (2 points). If I take both candies and the coin, that's 10 points. The remaining items would be 2 plates and 2 lemons. The other player might want those, but since I don't value them, it's okay if they take them. But I need to check if the other player's proposal would conflict with mine. If they also want the coin or candies, there's a problem. But since I'm taking the candies and coin, they can't take those. So their proposal should only include the plates and lemons. } \\
\\ 
\texttt{Wait, but the other player's value function is secret. Maybe they value the plates or lemons more. If they take the plates and lemons, that's fine. But if they also want the coin, then our proposals would conflict. But since I'm already taking the coin, they can't take it. So maybe the other player's proposal would be to take the plates and lemons. That way, our proposals are complementary. } \\
\\ 
\texttt{Alternatively, maybe I should leave the coin for the other player. If I take both candies (8 points) and leave the coin, plates, and lemons for them. Then their proposal could include the coin and maybe some plates or lemons. But if they take the coin, that's 2 points for them, and I have 8. That's a good split. } \\
\\ 
\texttt{But I need to make sure that the other player's proposal doesn't require items I've already taken. So if I take the candies and the coin, they can't take those. If they take the remaining items (plates and lemons), that's okay. But if their value function includes the coin, they might want it, but since I took it, they can't. So maybe the other player would be okay with taking the plates and lemons if they} \\
            }
        }
    }
    & & \\ \\

    \theutterance \stepcounter{utterance}  
    & & \multicolumn{4}{p{0.6\linewidth}}{
        \cellcolor[rgb]{0.9,0.9,0.9}{
            \makecell[{{p{\linewidth}}}]{
                \texttt{\tiny{[GM$\rangle$P2]}}
                \texttt{<think>} \\
\texttt{Okay, let's try to figure out the best way to approach this. My value function is plate: 0, coin: 2, candy: 4, lemon: 0. So the candies are the most valuable to me, followed by the coin. The other player's values are secret, so I don't know what they want. But I need to make sure that our proposals are complementary so we both get points.} \\
\\ 
\texttt{First, I should aim to get the items I value the most. That would be the two candies (4 each, so 8 points) and the coin (2 points). If I take both candies and the coin, that's 10 points. The remaining items would be 2 plates and 2 lemons. The other player might want those, but since I don't value them, it's okay if they take them. But I need to check if the other player's proposal would conflict with mine. If they also want the coin or candies, there's a problem. But since I'm taking the candies and coin, they can't take those. So their proposal should only include the plates and lemons. } \\
\\ 
\texttt{Wait, but the other player's value function is secret. Maybe they value the plates or lemons more. If they take the plates and lemons, that's fine. But if they also want the coin, then our proposals would conflict. But since I'm already taking the coin, they can't take it. So maybe the other player's proposal would be to take the plates and lemons. That way, our proposals are complementary. } \\
\\ 
\texttt{Alternatively, maybe I should leave the coin for the other player. If I take both candies (8 points) and leave the coin, plates, and lemons for them. Then their proposal could include the coin and maybe some plates or lemons. But if they take the coin, that's 2 points for them, and I have 8. That's a good split. } \\
\\ 
\texttt{But I need to make sure that the other player's proposal doesn't require items I've already taken. So if I take the candies and the coin, they can't take those. If they take the remaining items (plates and lemons), that's okay. But if their value function includes the coin, they might want it, but since I took it, they can't. So maybe the other player would be okay with taking the plates and lemons if they} \\
            }
        }
    }
    & \\ \\

    \theutterance \stepcounter{utterance}  
    & & & \multicolumn{4}{p{0.6\linewidth}}{
        \cellcolor[rgb]{0.9,0.9,0.9}{
            \makecell[{{p{\linewidth}}}]{
                \texttt{\tiny{[GM$\langle$P2]}}
                \texttt{<think>} \\
\texttt{Okay, let's try to figure out the best way to approach this. My value function is plate: 0, coin: 2, candy: 4, lemon: 0. So the candies are the most valuable to me, followed by the coin. The other player's values are secret, so I don't know what they want. But I need to make sure that our proposals are complementary so we both get points.} \\
\\ 
\texttt{First, I should aim to get the items I value the most. That would be the two candies (4 each, so 8 points) and the coin (2 points). If I take both candies and the coin, that's 10 points. The remaining items would be 2 plates and 2 lemons. The other player might want those, but since I don't value them, it's okay if they take them. But I need to check if the other player's proposal would conflict with mine. If they also want the coin or candies, there's a problem. But since I'm taking the candies and coin, they can't take those. So their proposal should only include the plates and lemons. } \\
\\ 
\texttt{Wait, but the other player's value function is secret. Maybe they value the plates or lemons more. If they take the plates and lemons, that's fine. But if they also want the coin, then our proposals would conflict. But since I'm already taking the coin, they can't take it. So maybe the other player's proposal would be to take the plates and lemons. That way, our proposals are complementary. } \\
\\ 
\texttt{Alternatively, maybe I should leave the coin for the other player. If I take both candies (8 points) and leave the coin, plates, and lemons for them. Then their proposal could include the coin and maybe some plates or lemons. But if they take the coin, that's 2 points for them, and I have 8. That's a good split. } \\
\\ 
\texttt{But I need to make sure that the other player's proposal doesn't require items I've already taken. So if I take the candies and the coin, they can't take those. If they take the remaining items (plates and lemons), that's okay. But if their value function includes the coin, they might want it, but since I took it, they can't. So maybe the other player would be okay with taking the plates and lemons if they} \\
            }
        }
    }
     \\ \\

    \theutterance \stepcounter{utterance}  
    & & & \multicolumn{2}{p{0.3\linewidth}}{
        \cellcolor[rgb]{0.9,0.9,0.9}{
            \makecell[{{p{\linewidth}}}]{
                \texttt{\tiny{[GM$|$GM]}}
                \texttt{continue} \\
            }
        }
    }
    & & \\ \\

    \theutterance \stepcounter{utterance}  
    & & & \multicolumn{2}{p{0.3\linewidth}}{
        \cellcolor[rgb]{0.9,0.9,0.9}{
            \makecell[{{p{\linewidth}}}]{
                \texttt{\tiny{[GM$|$GM]}}
                \texttt{<think>} \\
\texttt{Okay, let's try to figure out the best way to approach this. My value function is plate: 0, coin: 2, candy: 4, lemon: 0. So the candies are the most valuable to me, followed by the coin. The other player's values are secret, so I don't know what they want. But I need to make sure that our proposals are complementary so we both get points.} \\
\\ 
\texttt{First, I should aim to get the items I value the most. That would be the two candies (4 each, so 8 points) and the coin (2 points). If I take both candies and the coin, that's 10 points. The remaining items would be 2 plates and 2 lemons. The other player might want those, but since I don't value them, it's okay if they take them. But I need to check if the other player's proposal would conflict with mine. If they also want the coin or candies, there's a problem. But since I'm taking the candies and coin, they can't take those. So their proposal should only include the plates and lemons. } \\
\\ 
\texttt{Wait, but the other player's value function is secret. Maybe they value the plates or lemons more. If they take the plates and lemons, that's fine. But if they also want the coin, then our proposals would conflict. But since I'm already taking the coin, they can't take it. So maybe the other player's proposal would be to take the plates and lemons. That way, our proposals are complementary. } \\
\\ 
\texttt{Alternatively, maybe I should leave the coin for the other player. If I take both candies (8 points) and leave the coin, plates, and lemons for them. Then their proposal could include the coin and maybe some plates or lemons. But if they take the coin, that's 2 points for them, and I have 8. That's a good split. } \\
\\ 
\texttt{But I need to make sure that the other player's proposal doesn't require items I've already taken. So if I take the candies and the coin, they can't take those. If they take the remaining items (plates and lemons), that's okay. But if their value function includes the coin, they might want it, but since I took it, they can't. So maybe the other player would be okay with taking the plates and lemons if they} \\
            }
        }
    }
    & & \\ \\

    \theutterance \stepcounter{utterance}  
    & & \multicolumn{4}{p{0.6\linewidth}}{
        \cellcolor[rgb]{0.9,0.9,0.9}{
            \makecell[{{p{\linewidth}}}]{
                \texttt{\tiny{[P1$\langle$GM]}}
                \texttt{<think>} \\
\texttt{Okay, let's try to figure out the best way to approach this. My value function is plate: 0, coin: 2, candy: 4, lemon: 0. So the candies are the most valuable to me, followed by the coin. The other player's values are secret, so I don't know what they want. But I need to make sure that our proposals are complementary so we both get points.} \\
\\ 
\texttt{First, I should aim to get the items I value the most. That would be the two candies (4 each, so 8 points) and the coin (2 points). If I take both candies and the coin, that's 10 points. The remaining items would be 2 plates and 2 lemons. The other player might want those, but since I don't value them, it's okay if they take them. But I need to check if the other player's proposal would conflict with mine. If they also want the coin or candies, there's a problem. But since I'm taking the candies and coin, they can't take those. So their proposal should only include the plates and lemons. } \\
\\ 
\texttt{Wait, but the other player's value function is secret. Maybe they value the plates or lemons more. If they take the plates and lemons, that's fine. But if they also want the coin, then our proposals would conflict. But since I'm already taking the coin, they can't take it. So maybe the other player's proposal would be to take the plates and lemons. That way, our proposals are complementary. } \\
\\ 
\texttt{Alternatively, maybe I should leave the coin for the other player. If I take both candies (8 points) and leave the coin, plates, and lemons for them. Then their proposal could include the coin and maybe some plates or lemons. But if they take the coin, that's 2 points for them, and I have 8. That's a good split. } \\
\\ 
\texttt{But I need to make sure that the other player's proposal doesn't require items I've already taken. So if I take the candies and the coin, they can't take those. If they take the remaining items (plates and lemons), that's okay. But if their value function includes the coin, they might want it, but since I took it, they can't. So maybe the other player would be okay with taking the plates and lemons if they} \\
            }
        }
    }
    & \\ \\

    \theutterance \stepcounter{utterance}  
    & \multicolumn{4}{p{0.6\linewidth}}{
        \cellcolor[rgb]{0.9,0.9,0.9}{
            \makecell[{{p{\linewidth}}}]{
                \texttt{\tiny{[P1$\rangle$GM]}}
                \texttt{<think>} \\
\texttt{Okay, let's try to figure out the best way to approach this. My value function is plate: 0, coin: 2, candy: 4, lemon: 0. So the candies are the most valuable to me, followed by the coin. The other player's values are secret, so I don't know what they want. But I need to make sure that our proposals are complementary so we both get points.} \\
\\ 
\texttt{First, I should aim to get the items I value the most. That would be the two candies (4 each, so 8 points) and the coin (2 points). If I take both candies and the coin, that's 10 points. The remaining items would be 2 plates and 2 lemons. The other player might want those, but since I don't value them, it's okay if they take them. But I need to check if the other player's proposal would conflict with mine. If they also want the coin or candies, there's a problem. But since I'm taking the candies and coin, they can't take those. So their proposal should only include the plates and lemons. } \\
\\ 
\texttt{Wait, but the other player's value function is secret. Maybe they value the plates or lemons more. If they take the plates and lemons, that's fine. But if they also want the coin, then our proposals would conflict. But since I'm already taking the coin, they can't take it. So maybe the other player's proposal would be to take the plates and lemons. That way, our proposals are complementary. } \\
\\ 
\texttt{Alternatively, maybe I should leave the coin for the other player. If I take both candies (8 points) and leave the coin, plates, and lemons for them. Then their proposal could include the coin and maybe some plates or lemons. But if they take the coin, that's 2 points for them, and I have 8. That's a good split. } \\
\\ 
\texttt{But I need to make sure that the other player's proposal doesn't require items I've already taken. So if I take the candies and the coin, they can't take those. If they take the remaining items (plates and lemons), that's okay. But if their value function includes the coin, they might want it, but since I took it, they can't. So maybe the other player would be okay with taking the plates and lemons if they} \\
            }
        }
    }
    & & \\ \\

    \theutterance \stepcounter{utterance}  
    & & & \multicolumn{2}{p{0.3\linewidth}}{
        \cellcolor[rgb]{0.9,0.9,0.9}{
            \makecell[{{p{\linewidth}}}]{
                \texttt{\tiny{[GM$|$GM]}}
                \texttt{continue} \\
            }
        }
    }
    & & \\ \\

    \theutterance \stepcounter{utterance}  
    & & & \multicolumn{2}{p{0.3\linewidth}}{
        \cellcolor[rgb]{0.9,0.9,0.9}{
            \makecell[{{p{\linewidth}}}]{
                \texttt{\tiny{[GM$|$GM]}}
                \texttt{<think>} \\
\texttt{Okay, let's try to figure out the best way to approach this. My value function is plate: 0, coin: 2, candy: 4, lemon: 0. So the candies are the most valuable to me, followed by the coin. The other player's values are secret, so I don't know what they want. But I need to make sure that our proposals are complementary so we both get points.} \\
\\ 
\texttt{First, I should aim to get the items I value the most. That would be the two candies (4 each, so 8 points) and the coin (2 points). If I take both candies and the coin, that's 10 points. The remaining items would be 2 plates and 2 lemons. The other player might want those, but since I don't value them, it's okay if they take them. But I need to check if the other player's proposal would conflict with mine. If they also want the coin or candies, there's a problem. But since I'm taking the candies and coin, they can't take those. So their proposal should only include the plates and lemons. } \\
\\ 
\texttt{Wait, but the other player's value function is secret. Maybe they value the plates or lemons more. If they take the plates and lemons, that's fine. But if they also want the coin, then our proposals would conflict. But since I'm already taking the coin, they can't take it. So maybe the other player's proposal would be to take the plates and lemons. That way, our proposals are complementary. } \\
\\ 
\texttt{Alternatively, maybe I should leave the coin for the other player. If I take both candies (8 points) and leave the coin, plates, and lemons for them. Then their proposal could include the coin and maybe some plates or lemons. But if they take the coin, that's 2 points for them, and I have 8. That's a good split. } \\
\\ 
\texttt{But I need to make sure that the other player's proposal doesn't require items I've already taken. So if I take the candies and the coin, they can't take those. If they take the remaining items (plates and lemons), that's okay. But if their value function includes the coin, they might want it, but since I took it, they can't. So maybe the other player would be okay with taking the plates and lemons if they} \\
            }
        }
    }
    & & \\ \\

    \theutterance \stepcounter{utterance}  
    & & \multicolumn{4}{p{0.6\linewidth}}{
        \cellcolor[rgb]{0.9,0.9,0.9}{
            \makecell[{{p{\linewidth}}}]{
                \texttt{\tiny{[GM$\rangle$P2]}}
                \texttt{<think>} \\
\texttt{Okay, let's try to figure out the best way to approach this. My value function is plate: 0, coin: 2, candy: 4, lemon: 0. So the candies are the most valuable to me, followed by the coin. The other player's values are secret, so I don't know what they want. But I need to make sure that our proposals are complementary so we both get points.} \\
\\ 
\texttt{First, I should aim to get the items I value the most. That would be the two candies (4 each, so 8 points) and the coin (2 points). If I take both candies and the coin, that's 10 points. The remaining items would be 2 plates and 2 lemons. The other player might want those, but since I don't value them, it's okay if they take them. But I need to check if the other player's proposal would conflict with mine. If they also want the coin or candies, there's a problem. But since I'm taking the candies and coin, they can't take those. So their proposal should only include the plates and lemons. } \\
\\ 
\texttt{Wait, but the other player's value function is secret. Maybe they value the plates or lemons more. If they take the plates and lemons, that's fine. But if they also want the coin, then our proposals would conflict. But since I'm already taking the coin, they can't take it. So maybe the other player's proposal would be to take the plates and lemons. That way, our proposals are complementary. } \\
\\ 
\texttt{Alternatively, maybe I should leave the coin for the other player. If I take both candies (8 points) and leave the coin, plates, and lemons for them. Then their proposal could include the coin and maybe some plates or lemons. But if they take the coin, that's 2 points for them, and I have 8. That's a good split. } \\
\\ 
\texttt{But I need to make sure that the other player's proposal doesn't require items I've already taken. So if I take the candies and the coin, they can't take those. If they take the remaining items (plates and lemons), that's okay. But if their value function includes the coin, they might want it, but since I took it, they can't. So maybe the other player would be okay with taking the plates and lemons if they} \\
            }
        }
    }
    & \\ \\

    \theutterance \stepcounter{utterance}  
    & & & \multicolumn{4}{p{0.6\linewidth}}{
        \cellcolor[rgb]{0.9,0.9,0.9}{
            \makecell[{{p{\linewidth}}}]{
                \texttt{\tiny{[GM$\langle$P2]}}
                \texttt{<think>} \\
\texttt{Okay, let's try to figure out the best way to approach this. My value function is plate: 0, coin: 2, candy: 4, lemon: 0. So the candies are the most valuable to me, followed by the coin. The other player's values are secret, so I don't know what they want. But I need to make sure that our proposals are complementary so we both get points.} \\
\\ 
\texttt{First, I should aim to get the items I value the most. That would be the two candies (4 each, so 8 points) and the coin (2 points). If I take both candies and the coin, that's 10 points. The remaining items would be 2 plates and 2 lemons. The other player might want those, but since I don't value them, it's okay if they take them. But I need to check if the other player's proposal would conflict with mine. If they also want the coin or candies, there's a problem. But since I'm taking the candies and coin, they can't take those. So their proposal should only include the plates and lemons. } \\
\\ 
\texttt{Wait, but the other player's value function is secret. Maybe they value the plates or lemons more. If they take the plates and lemons, that's fine. But if they also want the coin, then our proposals would conflict. But since I'm already taking the coin, they can't take it. So maybe the other player's proposal would be to take the plates and lemons. That way, our proposals are complementary. } \\
\\ 
\texttt{Alternatively, maybe I should leave the coin for the other player. If I take both candies (8 points) and leave the coin, plates, and lemons for them. Then their proposal could include the coin and maybe some plates or lemons. But if they take the coin, that's 2 points for them, and I have 8. That's a good split. } \\
\\ 
\texttt{But I need to make sure that the other player's proposal doesn't require items I've already taken. So if I take the candies and the coin, they can't take those. If they take the remaining items (plates and lemons), that's okay. But if their value function includes the coin, they might want it, but since I took it, they can't. So maybe the other player would be okay with taking the plates and lemons if they} \\
            }
        }
    }
     \\ \\

    \theutterance \stepcounter{utterance}  
    & & & \multicolumn{2}{p{0.3\linewidth}}{
        \cellcolor[rgb]{0.9,0.9,0.9}{
            \makecell[{{p{\linewidth}}}]{
                \texttt{\tiny{[GM$|$GM]}}
                \texttt{continue} \\
            }
        }
    }
    & & \\ \\

    \theutterance \stepcounter{utterance}  
    & & & \multicolumn{2}{p{0.3\linewidth}}{
        \cellcolor[rgb]{0.9,0.9,0.9}{
            \makecell[{{p{\linewidth}}}]{
                \texttt{\tiny{[GM$|$GM]}}
                \texttt{<think>} \\
\texttt{Okay, let's try to figure out the best way to approach this. My value function is plate: 0, coin: 2, candy: 4, lemon: 0. So the candies are the most valuable to me, followed by the coin. The other player's values are secret, so I don't know what they want. But I need to make sure that our proposals are complementary so we both get points.} \\
\\ 
\texttt{First, I should aim to get the items I value the most. That would be the two candies (4 each, so 8 points) and the coin (2 points). If I take both candies and the coin, that's 10 points. The remaining items would be 2 plates and 2 lemons. The other player might want those, but since I don't value them, it's okay if they take them. But I need to check if the other player's proposal would conflict with mine. If they also want the coin or candies, there's a problem. But since I'm taking the candies and coin, they can't take those. So their proposal should only include the plates and lemons. } \\
\\ 
\texttt{Wait, but the other player's value function is secret. Maybe they value the plates or lemons more. If they take the plates and lemons, that's fine. But if they also want the coin, then our proposals would conflict. But since I'm already taking the coin, they can't take it. So maybe the other player's proposal would be to take the plates and lemons. That way, our proposals are complementary. } \\
\\ 
\texttt{Alternatively, maybe I should leave the coin for the other player. If I take both candies (8 points) and leave the coin, plates, and lemons for them. Then their proposal could include the coin and maybe some plates or lemons. But if they take the coin, that's 2 points for them, and I have 8. That's a good split. } \\
\\ 
\texttt{But I need to make sure that the other player's proposal doesn't require items I've already taken. So if I take the candies and the coin, they can't take those. If they take the remaining items (plates and lemons), that's okay. But if their value function includes the coin, they might want it, but since I took it, they can't. So maybe the other player would be okay with taking the plates and lemons if they} \\
            }
        }
    }
    & & \\ \\

    \theutterance \stepcounter{utterance}  
    & & \multicolumn{4}{p{0.6\linewidth}}{
        \cellcolor[rgb]{0.9,0.9,0.9}{
            \makecell[{{p{\linewidth}}}]{
                \texttt{\tiny{[P1$\langle$GM]}}
                \texttt{<think>} \\
\texttt{Okay, let's try to figure out the best way to approach this. My value function is plate: 0, coin: 2, candy: 4, lemon: 0. So the candies are the most valuable to me, followed by the coin. The other player's values are secret, so I don't know what they want. But I need to make sure that our proposals are complementary so we both get points.} \\
\\ 
\texttt{First, I should aim to get the items I value the most. That would be the two candies (4 each, so 8 points) and the coin (2 points). If I take both candies and the coin, that's 10 points. The remaining items would be 2 plates and 2 lemons. The other player might want those, but since I don't value them, it's okay if they take them. But I need to check if the other player's proposal would conflict with mine. If they also want the coin or candies, there's a problem. But since I'm taking the candies and coin, they can't take those. So their proposal should only include the plates and lemons. } \\
\\ 
\texttt{Wait, but the other player's value function is secret. Maybe they value the plates or lemons more. If they take the plates and lemons, that's fine. But if they also want the coin, then our proposals would conflict. But since I'm already taking the coin, they can't take it. So maybe the other player's proposal would be to take the plates and lemons. That way, our proposals are complementary. } \\
\\ 
\texttt{Alternatively, maybe I should leave the coin for the other player. If I take both candies (8 points) and leave the coin, plates, and lemons for them. Then their proposal could include the coin and maybe some plates or lemons. But if they take the coin, that's 2 points for them, and I have 8. That's a good split. } \\
\\ 
\texttt{But I need to make sure that the other player's proposal doesn't require items I've already taken. So if I take the candies and the coin, they can't take those. If they take the remaining items (plates and lemons), that's okay. But if their value function includes the coin, they might want it, but since I took it, they can't. So maybe the other player would be okay with taking the plates and lemons if they} \\
\\ 
\\ 
\texttt{The time is up. It is now your turn to submit a secret proposal.} \\
            }
        }
    }
    & \\ \\

    \theutterance \stepcounter{utterance}  
    & \multicolumn{4}{p{0.6\linewidth}}{
        \cellcolor[rgb]{0.9,0.9,0.9}{
            \makecell[{{p{\linewidth}}}]{
                \texttt{\tiny{[P1$\rangle$GM]}}
                \texttt{<think>} \\
\texttt{Okay, let's try to figure out the best way to approach this. My value function is plate: 0, coin: 2, candy: 4, lemon: 0. So the candies are the most valuable to me, followed by the coin. The other player's values are secret, so I don't know what they want. But I need to make sure that our proposals are complementary so we both get points.} \\
\\ 
\texttt{First, I should aim to get the items I value the most. That would be the two candies (4 each, so 8 points) and the coin (2 points). If I take both candies and the coin, that's 10 points. The remaining items would be 2 plates and 2 lemons. The other player might want those, but since I don't value them, it's okay if they take them. But I need to check if the other player's proposal would conflict with mine. If they also want the coin or candies, there's a problem. But since I'm taking the candies and coin, they can't take those. So their proposal should only include the plates and lemons. } \\
\\ 
\texttt{Wait, but the other player's value function is secret. Maybe they value the plates or lemons more. If they take the plates and lemons, that's fine. But if they also want the coin, then our proposals would conflict. But since I'm already taking the coin, they can't take it. So maybe the other player's proposal would be to take the plates and lemons. That way, our proposals are complementary. } \\
\\ 
\texttt{Alternatively, maybe I should leave the coin for the other player. If I take both candies (8 points) and leave the coin, plates, and lemons for them. Then their proposal could include the coin and maybe some plates or lemons. But if they take the coin, that's 2 points for them, and I have 8. That's a good split. } \\
\\ 
\texttt{But I need to make sure that the other player's proposal doesn't require items I've already taken. So if I take the candies and the coin, they can't take those. If they take the remaining items (plates and lemons), that's okay. But if their value function includes the coin, they might want it, but since I took it, they can't. So maybe the other player would be okay with taking the plates and lemons if they} \\
            }
        }
    }
    & & \\ \\

    \theutterance \stepcounter{utterance}  
    & & & \multicolumn{2}{p{0.3\linewidth}}{
        \cellcolor[rgb]{0.9,0.9,0.9}{
            \makecell[{{p{\linewidth}}}]{
                \texttt{\tiny{[GM$|$GM]}}
                \texttt{continue} \\
            }
        }
    }
    & & \\ \\

    \theutterance \stepcounter{utterance}  
    & & & \multicolumn{2}{p{0.3\linewidth}}{
        \cellcolor[rgb]{0.9,0.9,0.9}{
            \makecell[{{p{\linewidth}}}]{
                \texttt{\tiny{[GM$|$GM]}}
                \texttt{player was instructed to make a proposal} \\
            }
        }
    }
    & & \\ \\

\end{supertabular}
}

\end{document}
